\section{Methodology}\label{sec:method}

\subsection{Notation and problem setup}

Let $y_t$ denote the NEMA hourly load at time $t$. Forecasting is performed from a
fixed-length lookback window of $L = 168$ hours, i.e. the one-week interval
$[t-L,\, t-1]$. For each forecast origin $t$ and each horizon
$h \in \{1, \dots, 24\}$, the task is to predict the load $y_{t+h}$ at the target
hour $t+h$. We adopt a \emph{direct} multi-horizon strategy: a separate model is
trained for each horizon $h$, mapping a feature vector to a single scalar target
$y_{t+h}$, with no recursive feedback of intermediate predictions.

\subsection{Window lag features}

The first feature block summarizes the lookback window
$[t-L,\, t-1]$. From each engineered channel we extract values at lags
$\{1, 4, 8, 24, 48, 168\}$ hours and rolling means over windows of $\{24, 168\}$
hours, capturing recent dynamics together with daily and weekly periodicity. The
raw load channel (RTLO) is restricted to a whitelist retained from a notebook
ablation --- \texttt{RTLO\_lag4}, \texttt{RTLO\_mean24}, \texttt{RTLO\_lag48},
\texttt{RTLO\_lag168}, and \texttt{RTLO\_mean168} --- which avoids redundant
near-term load lags while preserving the daily and weekly
load signal. The full window feature map, denoted $\phi_{\mathrm{lag}}(\cdot)$,
produces 173 features.

\subsection{Calendar features}

Calendar effects are encoded with cyclic (sine/cosine) transforms of the hour of
day, day of week, and month, together with binary indicators \texttt{is\_weekend}
and \texttt{is\_us\_holiday}. The cyclic encoding avoids artificial discontinuities
at period boundaries (e.g. hour 23 to hour 0) and lets tree splits exploit smooth
diurnal, weekly, and seasonal structure.

\subsection{Weather-derived features}

Temperature drives electricity demand non-linearly through both heating and
cooling load. We compute heating and cooling degree days relative to a 65$^{\circ}$F
base,
\begin{equation}
\mathrm{HDD} = \max(0,\ 65 - T), \qquad
\mathrm{CDD} = \max(0,\ T - 65),
\end{equation}
together with the squared terms $T^2$, $\mathrm{HDD}^2$, and $\mathrm{CDD}^2$ to
capture curvature, and sigmoidal thermal-stress transforms (logistic functions
centred at 18$^{\circ}$C for heating and 22$^{\circ}$C for cooling) that model the
saturating onset of heating and air-conditioning demand. These derived quantities
are computed from the Open-Meteo weather variables described in
Section~\ref{sec:data}.

\subsection{Target-hour exogenous features}

The central methodological idea is to provide each horizon model with information
about the target hour $t+h$ itself, rather than only the history up to $t-1$. The
target-hour exogenous block $\phi_{\mathrm{exog}}(t+h)$ comprises all engineered
calendar features and all weather-derived features \emph{evaluated at the target
hour} $t+h$ --- the exactly-known calendar attributes of $t+h$ and the
\emph{forecasted} weather at $t+h$. The raw load channel RTLO is excluded from
this block, since the load at $t+h$ is precisely the prediction target. This block
contributes 21 features.

\subsection{Augmented feature map}

For each horizon $h$, the window features and the target-hour exogenous features
are concatenated into a single augmented vector:
\begin{equation}
\mathbf{x}_h \;=\; \bigl[\, \phi_{\mathrm{lag}}\!\bigl([t-L,\, t-1]\bigr)
\;;\; \phi_{\mathrm{exog}}(t+h) \,\bigr],
\label{eq:augmented}
\end{equation}
yielding a $173 + 21 = 194$-dimensional representation. Each horizon model
$f_h$ then predicts
\begin{equation}
\hat{y}_{t+h} \;=\; f_h(\mathbf{x}_h).
\end{equation}
The window block is shared in form across horizons, while the exogenous block
$\phi_{\mathrm{exog}}(t+h)$ shifts to the appropriate target hour for each $h$.

\subsection{Direct multi-horizon strategy}

We instantiate 24 independent CatBoost models, $\{f_h\}_{h=1}^{24}$, one per
horizon, each trained directly on its own target $y_{t+h}$. This contrasts with a
recursive roll-out, in which a single one-step model is applied iteratively, its
predictions fed back as inputs for subsequent steps. Recursive roll-out compounds
errors across steps and suffers from a train/serve mismatch, because at training
time the model conditions on true past values whereas at inference time it
conditions on its own noisy forecasts~\cite{taieb2016bias}. The direct strategy
trains each horizon under exactly the conditions it faces at deployment, removing
both error compounding and that mismatch, at the cost of fitting 24 models. For
ablation we also consider a single $h{=}1$ model whose prediction is reused for
every horizon.

\subsection{Learner and training configuration}

Each $f_h$ is a CatBoost gradient-boosted tree ensemble~\cite{prokhorenkova2018catboost},
chosen for strong tabular performance and native handling of the heterogeneous
engineered features. All 24 models share the same configuration: 1000 boosting
iterations, tree depth 8, learning rate 0.05, mean-absolute-error (MAE) loss,
early stopping after 50 rounds without validation improvement, and random seed 42.
The MAE objective aligns training directly with the primary evaluation metric and
is robust to occasional load spikes. Early stopping uses the validation split
defined in Section~\ref{sec:data}.

\subsection{Why target-hour weather is causal, not leakage}

A natural concern is whether conditioning on quantities at $t+h$ introduces
target leakage. It does not. The target-hour weather in $\phi_{\mathrm{exog}}(t+h)$
is a \emph{forecast}, which is exactly the quantity available to an operator at
prediction time; in deployment it is supplied by the Open-Meteo forecast API, and
in evaluation the corresponding archive value plays this role. The calendar
attributes of $t+h$ are deterministic and known arbitrarily far in advance.
Crucially, the load channel RTLO at $t+h$ --- the only quantity that would
constitute leakage --- is explicitly excluded from the exogenous block, and every
lag and rolling feature in $\phi_{\mathrm{lag}}$ references data at or before
$t-1$. No feature can therefore carry target-time load information, a property we
verify empirically with a shuffled-target audit in the results.
